\subsection{Pages 7--8: Deriving the Gradient}

\subsubsection*{Page 7 --- Gradient of the RL Objective}

\textbf{Quick recap.}
Remember from page 6: the agent's goal is to maximize $J(\theta)$, the average total reward it
gets by following its policy. To improve the policy, we need to compute the \textbf{gradient}
of $J(\theta)$ --- basically, the direction in which to nudge the policy knobs ($\theta$) to
get more reward.

\medskip
\textbf{Equation 1 --- Rewriting the objective as an integral.}

\[
  J(\theta)
  = E_{\tau \sim p_\theta(\tau)}\!\left[\sum_t r(s_t, a_t)\right]
  = \int p_\theta(\tau)\,r(\tau)\,d\tau
\]

\begin{tabular}{ll}
  Symbol & Meaning \\[4pt]
  $J(\theta)$ & The score of the policy --- total expected reward \\
  $\theta$ & The policy's knobs (parameters) \\
  $\tau$ & A trajectory --- a full playthrough: $s_1,a_1,s_2,a_2,\ldots$ \\
  $p_\theta(\tau)$ & Probability of that trajectory happening under policy $\theta$ \\
  $r(\tau)$ & Total reward collected along that trajectory \\
  $\int\ldots\,d\tau$ & ``Sum over all possible trajectories'' \\
\end{tabular}

\medskip
This is just rewriting the expected reward in math notation. Instead of saying ``average reward
over many runs,'' we write it as an integral --- summing $p_\theta(\tau)\cdot r(\tau)$ over
every possible trajectory. The integral is mathematically equivalent to the expectation, which
sets us up for taking the gradient.

\medskip
\textbf{Equation 2 --- Taking the gradient (the log-derivative trick).}

\[
  \nabla_\theta J(\theta)
  = \int \nabla_\theta p_\theta(\tau)\,r(\tau)\,d\tau
  = \int p_\theta(\tau)\,\nabla_\theta \log p_\theta(\tau)\,r(\tau)\,d\tau
  = E_{\tau \sim p_\theta}\!\left[\nabla_\theta \log p_\theta(\tau)\,r(\tau)\right]
\]

\begin{tabular}{ll}
  Symbol & Meaning \\[4pt]
  $\nabla_\theta$ & ``How much does this change when we tweak $\theta$?'' (the gradient) \\
  $\nabla_\theta p_\theta(\tau)$ & Gradient of the trajectory probability --- hard to work with directly \\
  $\nabla_\theta \log p_\theta(\tau)$ & Gradient of the \textit{log} of the trajectory probability --- much easier \\
\end{tabular}

\medskip
We want the gradient of $J(\theta)$, but differentiating $p_\theta(\tau)$ directly is messy
because it involves the unknown environment dynamics. The \textbf{log-derivative trick} says:
$p_\theta(\tau)\cdot\nabla_\theta\log p_\theta(\tau) = \nabla_\theta p_\theta(\tau)$.
This lets us swap the hard term out and rewrite the whole gradient as an \textbf{expectation} ---
which means we can estimate it just by running the policy and averaging, no integrals needed.

\medskip
\textbf{The key identity.}

\[
  p_\theta(\tau)\,\nabla_\theta \log p_\theta(\tau)
  = p_\theta(\tau)\cdot\frac{\nabla_\theta p_\theta(\tau)}{p_\theta(\tau)}
  = \nabla_\theta p_\theta(\tau)
\]

This is just the chain rule of calculus: the derivative of $\log(x)$ is $\frac{1}{x}$.
Multiplying $p_\theta(\tau)$ on both sides cancels it out, proving the substitution above is valid.

\subsubsection*{Page 8 --- Expanding the Gradient into Something Computable}

Now the slide answers: ``Okay, but how do we actually \textbf{compute}
$\nabla_\theta\log p_\theta(\tau)$?''

\medskip
\textbf{Equation 3 --- What is $p_\theta(\tau)$?}

\[
  p_\theta(\tau)
  = p(s_1)\prod_{t=1}^{T}\pi_\theta(a_t\mid s_t)\,p(s_{t+1}\mid s_t,a_t)
\]

\begin{tabular}{ll}
  Symbol & Meaning \\[4pt]
  $p(s_1)$ & Probability of starting in state $s_1$ (set by the environment, not the agent) \\
  $\prod_{t=1}^T$ & Multiply together across all time steps $t=1$ to $T$ \\
  $\pi_\theta(a_t\mid s_t)$ & Policy's probability of picking action $a_t$ when in state $s_t$ \\
  $p(s_{t+1}\mid s_t,a_t)$ & Environment's probability of moving to $s_{t+1}$ given state and action \\
\end{tabular}

\medskip
A trajectory's probability is the product of every decision and every world transition along
the way. Think of it like the probability of a specific path through a maze --- it's the
starting probability times the chance you took each turn times the chance the world responded
the way it did, multiplied across every step.

\medskip
\textbf{Equation 4 --- Expanding $\log p_\theta(\tau)$.}

\[
  \nabla_\theta\!\left[
    \log p(s_1)
    + \sum_t \log\pi_\theta(a_t\mid s_t)
    + \log p(s_{t+1}\mid s_t,a_t)
  \right]
\]

Taking the log of a product turns it into a sum (since $\log(ABC)=\log A+\log B+\log C$).
The critical insight is that $\log p(s_1)$ and $\log p(s_{t+1}\mid s_t,a_t)$ do not contain
$\theta$ at all --- they are controlled by the environment, not the policy. So their gradients
are \textbf{zero} and they vanish. Only $\log\pi_\theta(a_t\mid s_t)$ survives.

\medskip
\textbf{Equation 5 --- The practical gradient formula.}

\[
  \nabla_\theta J(\theta)
  = E_{\tau\sim p_\theta}\!\left[
      \left(\sum_{t=1}^T \nabla_\theta\log\pi_\theta(a_t\mid s_t)\right)
      \left(\sum_{t=1}^T r(s_t,a_t)\right)
    \right]
\]

\begin{tabular}{ll}
  Symbol & Meaning \\[4pt]
  $\sum_{t=1}^T\nabla_\theta\log\pi_\theta(a_t\mid s_t)$ & How to nudge the policy to make each action more or less likely \\
  $\sum_{t=1}^T r(s_t,a_t)$ & Total reward collected during the trajectory \\
\end{tabular}

\medskip
The gradient is the product of two sums --- ``how much did each action's probability change
when we tweak $\theta$?'' multiplied by ``how much total reward did we get?'' If a trajectory
got high reward, we push $\theta$ in the direction that made those actions more likely. If it
got low reward, we push the other way. Crucially, the environment dynamics
$p(s_{t+1}\mid s_t,a_t)$ are \textbf{completely gone} --- we only need the policy and
the rewards.

\medskip
\textbf{Equation 6 --- The sample estimate (how we actually compute it).}

\[
  \nabla_\theta J(\theta)
  \approx \frac{1}{N}\sum_i
    \left(\sum_{t=1}^T\nabla_\theta\log\pi_\theta(a_{i,t}\mid s_{i,t})\right)
    \left(\sum_{t=1}^T r(s_{i,t},a_{i,t})\right)
\]

\begin{tabular}{ll}
  Symbol & Meaning \\[4pt]
  $N$ & Number of sampled trajectories (rollouts) \\
  $i$ & Index for each sampled trajectory \\
  $\frac{1}{N}\sum_i$ & Average over all $N$ rollouts \\
  $a_{i,t},\ s_{i,t}$ & Action and state at time $t$ in the $i$-th trajectory \\
\end{tabular}

\medskip
Since we cannot compute an expectation over \textit{all} possible trajectories, we approximate
it by running the policy $N$ times, computing the gradient formula for each run, and averaging.
This is the same idea as estimating the average of a die roll by rolling it many times --- the
more rollouts, the better the estimate. This is the actual algorithm you would implement in
code.


\subsection{Page 9: The Full REINFORCE Algorithm}

This page puts everything from pages 7 and 8 together into one complete, runnable algorithm.
It has a name: \textbf{REINFORCE}, also called the \textbf{vanilla policy gradient}.
``Vanilla'' just means the plain, basic version --- no fancy extras yet.

Think of it like a recipe with 4 steps that you repeat over and over until the agent gets good.

\medskip
\textbf{Step 1 --- Collect experience.}

\[
  \text{Sample } \{\tau^i\} \text{ from } \pi_\theta(a_t\mid s_t)
\]

\begin{tabular}{ll}
  Symbol & Meaning \\[4pt]
  $\{\tau^i\}$ & A collection of trajectories (multiple full playthroughs), indexed by $i$ \\
  $\pi_\theta(a_t\mid s_t)$ & The current policy --- given state $s_t$, it picks action $a_t$ \\
\end{tabular}

\medskip
Run your current policy in the environment multiple times and record what happened each time ---
which states you visited, which actions you took, and what rewards you got. This is just the
agent ``playing the game'' several times and taking notes.

\medskip
\textbf{Step 2 --- Compute the gradient.}

\[
  \nabla_\theta J(\theta)
  \approx \sum_i
    \left(\sum_{t=1}^T\nabla_\theta\log\pi_\theta(a_t^i\mid s_t^i)\right)
    \left(\sum_{t=1}^T r(s_t^i,a_t^i)\right)
\]

\begin{tabular}{ll}
  Symbol & Meaning \\[4pt]
  $\nabla_\theta J(\theta)$ & The gradient --- which direction to push the policy knobs \\
  $\sum_i$ & Sum (add up) over all collected trajectories \\
  $\nabla_\theta\log\pi_\theta(a_t^i\mid s_t^i)$ & How much action $a_t$ in trajectory $i$ would change if we tweaked $\theta$ \\
  $\sum_{t=1}^T r(s_t^i,a_t^i)$ & Total reward collected in trajectory $i$ \\
\end{tabular}

\medskip
For each recorded trajectory, multiply ``how sensitive was the policy to its own choices'' by
``how much reward did it get?'' Add this up across all trajectories. Trajectories with high
reward push the gradient hard in a direction that makes those actions more likely; trajectories
with low reward push in the opposite direction. This is the signal that tells us how to improve.

\medskip
\textbf{Step 3 --- Update the policy.}

\[
  \vec{\theta} = \vec{\theta} + \alpha\,\nabla_\theta J(\theta)
\]

\begin{tabular}{ll}
  Symbol & Meaning \\[4pt]
  $\vec{\theta}$ & The policy's parameters (knobs) --- a vector of numbers \\
  $\alpha$ & Learning rate --- controls how big a step to take \\
  $\nabla_\theta J(\theta)$ & The gradient computed in Step 2 \\
\end{tabular}

\medskip
Nudge the policy parameters $\theta$ in the direction of the gradient by a small amount
controlled by $\alpha$. Note it is a \textbf{plus} sign, not minus --- because we are doing
\textbf{gradient ascent} (climbing toward higher reward), unlike regular gradient descent used
in supervised learning which minimizes a loss. A large $\alpha$ means big jumps (risky); a
small $\alpha$ means careful, slow improvement.

\medskip
\textbf{Step 4 --- Repeat until convergence.}

Loop back to Step~1 with the newly updated policy. Collect new data, compute a new gradient,
update again. Keep going until the policy stops improving noticeably --- that is convergence.

\medskip
\textbf{The big picture.} The diagram on the slide shows this loop: collect experience by
running the policy, then improve the policy using that data, then repeat. REINFORCE looks at
all collected trajectories, figures out which actions led to good outcomes, and makes those
actions more likely next time. It is literally trial and error, formalized into math.


\subsection{Page 10: Understanding What the Gradient Actually Does}

This page is not introducing new math. It is stepping back and asking: \textbf{what does this
gradient formula actually mean in plain English?} It is the intuition behind everything derived
on pages 7--9.

\medskip
\textbf{The gradient (same formula, new lens).}

\[
  \nabla_\theta J(\theta)
  \approx \frac{1}{N}\sum_i
    \left(\sum_{t=1}^T\nabla_\theta\log\pi_\theta(a_t^i\mid s_t^i)\right)
    \left(\sum_{t=1}^T r(s_t^i,a_t^i)\right)
\]

\begin{tabular}{ll}
  Symbol & Meaning \\[4pt]
  $\frac{1}{N}\sum_i$ & Average over all $N$ collected trajectories \\
  $\nabla_\theta\log\pi_\theta(a_t^i\mid s_t^i)$ & Direction to push $\theta$ to make action $a_t$ in trajectory $i$ more likely \\
  $\sum_{t=1}^T r(s_t^i,a_t^i)$ & Total reward of trajectory $i$ --- the weight \\
\end{tabular}

\medskip
The total reward of a trajectory acts as a \textbf{multiplier} on the gradient. If a trajectory
earned high reward, its gradient contribution gets multiplied by a big positive number ---
meaning those actions get strongly encouraged. If the reward was negative, the multiplier
flips the sign and those actions get discouraged. The gradient is literally a weighted vote
from every trajectory.

\medskip
\textbf{The three key points from the slide.}

\begin{enumerate}
  \item \textbf{Increase the likelihood of actions you took in high-reward trajectories.}
        If a playthrough went well, assume the actions you took during it were probably good.
        Make those actions more likely in the future.

  \item \textbf{Decrease the likelihood of actions you took in negative-reward trajectories.}
        If a playthrough went badly, assume those actions were probably bad. Make them less
        likely going forward.

  \item \textbf{Do more of the good stuff, less of the bad stuff.}
        That is the whole algorithm in one sentence. The math just makes this idea precise
        and automatic.
\end{enumerate}

\textbf{The diagram.} The picture shows several trajectories starting from the same point.
Some end in failure and one ends in success. After running REINFORCE, the policy will shift
probability mass toward whatever sequence of actions led to the successful path, and away
from the actions that led to failure. Over many iterations, the successful path becomes more
and more likely.

\textbf{Why this page matters.} Everything on pages 7 and 8 was dense derivation. This page
is the payoff --- it tells you that all that math ultimately boils down to one dead-simple
idea: \textbf{reward good behavior, punish bad behavior, repeat.} REINFORCE is just that idea,
expressed precisely enough for a computer to execute.


\subsection{Pages 11--12: A Problem with the Gradient and the Reward-to-Go Fix}

\subsubsection*{Page 11 --- Quiz: A Problem With the Gradient}

This page gives a concrete worked example to expose a flaw hiding inside the gradient formula.
The task: \textbf{teach a robot to walk forward.}

The reward is simple:
\[
  r(s,a) = \text{forward velocity of the robot (negative if it goes backwards)}
\]

The agent runs 5 trajectories ($\tau^1$ through $\tau^5$):

\begin{tabular}{lll}
  Trajectory & What happened & Reward \\[4pt]
  $\tau^1$ & Falls backwards & Negative \\
  $\tau^2$ & Falls \textbf{forward} & Small positive \\
  $\tau^3$ & Stands still & Zero \\
  $\tau^4$ & One small step forward, then falls backwards & Mixed \\
  $\tau^5$ & One large step backwards, then small step forwards & Negative overall \\
\end{tabular}

\medskip
The question the slide asks: given these 5 trajectories and the gradient formula from page~9,
\textbf{what will the gradient actually encourage the policy to do?}

The answer: \textbf{it will encourage the robot to fall forward --- not to actually take a
step forward.}

\medskip
\textbf{Why? The flaw explained.}

$\tau^2$ (falls forward) gets a positive total reward because at least it was moving in the
right direction for a moment. $\tau^4$ (takes a step, then falls backward) might net out to a
small or zero total reward because the backward fall cancels the forward step.

So when the gradient formula multiplies each trajectory's actions by its \textbf{total reward},
$\tau^2$ gets a stronger positive push than $\tau^4$. The gradient is essentially saying:
``falling forward is better than stepping forward and then falling backward'' --- which is
wrong for the actual goal of learning to walk.

This is a \textbf{credit assignment problem}: the total reward of a whole trajectory is being
used to judge every single action inside it, even actions that happened early and had nothing
to do with the final outcome.

\subsubsection*{Page 12 --- Improving the Gradient (Reward-to-Go)}

\textbf{The key insight.}

\begin{quote}
  Policy behavior at time $t$ does not affect rewards at time $t' < t$.
\end{quote}

In plain English: \textbf{your action right now cannot change what already happened in the
past.} If you fall at step 5, that cannot hurt or help your reward from steps 1--4. Those
are already done. So it is unfair --- and noisy --- to judge an action at time $t$ by the
rewards that came \textbf{before} it. We should only hold an action responsible for the
rewards that come \textbf{after} it (including the current moment). These future rewards are
called the \textbf{reward-to-go}.

\medskip
\textbf{Old gradient (the flawed one).}

\[
  \nabla_\theta J(\theta)
  \approx \frac{1}{N}\sum_i
    \left(\sum_{t=1}^T\nabla_\theta\log\pi_\theta(a_t^i\mid s_t^i)\right)
    \left(\sum_{t=1}^T r(s_t^i,a_t^i)\right)
\]

The second sum is the \textbf{total reward of the entire trajectory} --- it judges every
action by the sum of all rewards, past and future.

\medskip
\textbf{New gradient (the fix).}

\[
  \nabla_\theta J(\theta)
  \approx \frac{1}{N}\sum_i\sum_{t=1}^T
    \nabla_\theta\log\pi_\theta(a_t^i\mid s_t^i)
    \left(\sum_{t'=t}^{T} r(s_{t'}^i,a_{t'}^i)\right)
\]

\begin{tabular}{ll}
  Symbol & Meaning \\[4pt]
  $\sum_{t=1}^T$ & Outer sum now goes over each individual timestep $t$, not just trajectories \\
  $\nabla_\theta\log\pi_\theta(a_t^i\mid s_t^i)$ & Same as before --- how to nudge $\theta$ for action $a_t$ \\
  $\sum_{t'=t}^{T} r(s_{t'}^i,a_{t'}^i)$ & \textbf{Reward-to-go}: rewards from time $t$ onward only \\
\end{tabular}

\medskip
Instead of judging action $a_t$ by the total reward of the whole episode, we only judge it by
the rewards it could have actually influenced --- everything from time $t$ to the end. This is
fairer and makes the gradient less noisy, because past rewards that had nothing to do with
action $a_t$ are no longer polluting its signal.

\medskip
\textbf{Back to the walking robot.} With reward-to-go, $\tau^4$ (step forward, then fall back)
gets properly evaluated. The forward step at time $t=1$ gets credited with the rewards from
$t=1$ onward --- which include the positive forward velocity from that step. The backward fall
at time $t=2$ only gets blamed for rewards from $t=2$ onward. The gradient now correctly
separates the good action from the bad one within the same trajectory, instead of blurring
them together under one total reward number.